% from main.tex
%%%%%%%%%%  scan idx 536  |  folio 523  %%%%%%%%%%
Rappelons (IV 2.10) que le bifoncteur
\[
(E,F)\rightsquigarrow E\otimes_{\mathcal O_X}F,
\]
où $E$ est un $\mathcal O_X$-module localement libre et $F$ un $\mathcal O_X$-module cohérent, induit une loi de composition
\begin{equation}\tag{1.3.1}
K^\circ(X)\times K_\bullet(X)\longrightarrow K_\bullet(X)
\end{equation}
qui fait de $K_\bullet(X)$ un $K^\circ(X)$-module unitaire. Le résultat-clé du présent exposé est le

\textbf{Théorème 1.3.2.} Soit $X$ un schéma noethérien. Pour tout $i,j\in\mathbb Z$ la loi de composition (1.3.1) induit par restriction une loi de composition
\begin{equation}\tag{1.3.3}
\Filt^i(X)\times\Filt_j(X)\longrightarrow\Filt_{j-i}(X).
\end{equation}

\textbf{Corollaire 1.3.3.} Si $X$ est un schéma noethérien de dimension $\leq j$, alors l'homomorphisme canonique
\[
\Theta_X:K^\circ(X)\longrightarrow K_\bullet(X)
\]
défini par $\Theta_X(x)=x\cdot\cl_\bullet(\mathcal O_X)$ ($x\in K^\circ(X)$) induit pour tout $i\in\mathbb Z$ un homomorphisme
\[
\Filt^i(X)\longrightarrow\Filt_{j-i}(X).
\]

\emph{Démonstration :} Montrons d'abord que (1.3.3), pour $i+j\leq k\in\mathbb Z$ et pour tout sous-schéma intègre fermé de dimension $\leq j$, implique (1.3.2) pour $i+j\leq k$. Grâce à (1.1.4) il suffit de montrer que pour tout $x\in\Filt^i(X)$ et pour tout sous-schéma intègre fermé $Z$ de dimension $\leq j$ on a
